\documentclass[10pt,a4paper,sans]{moderncv}
\moderncvstyle{classic}
\moderncvcolor{blue}
\usepackage[utf8]{inputenc}
\usepackage[scale=0.85]{geometry}
\usepackage{helvet}

\name{Alexandro}{Disla}
\title{Suivi, Évaluation, Redevabilité et Apprentissage (MEAL)}
\address{Rue Saint-Louis Jeanty \#14, Route de Frère}{Port-au-Prince, Haïti}{}
\phone[mobile]{(+509) 4148-3700}
\email{alexandrodisla@hotmail.com}
\extrainfo{GitHub: \url{https://github.com/AD0791}}

\begin{document}
\makecvtitle

\section{Profil Professionnel}
\cvitem{}{Praticien MEAL avec \textbf{cinq ans d'expérience} sur des projets multi-sites en Haïti, répartis sur trois secteurs --- santé et VIH à la Caris Foundation International, eau et assainissement chez HANWASH, éducation chez Anseye Pou Ayiti --- précédés de huit ans de suivi d'indicateurs de projets publics au Ministère de la Planification. Je conçois les outils de collecte sur \textbf{KoboToolbox} et \textbf{ODK}, forme les agents qui les utilisent, vérifie les données remontées, en tire les indicateurs et les restitue aux équipes. \textbf{Ingénieur de données} en parallèle, j'administre les bases et construis les tableaux de bord dont dépend un système MEAL. Résident de Port-au-Prince ; \textbf{français et créole haïtien maternels}, anglais courant.}

\section{Compétences Clés}
\cvitem{}{
\begin{itemize}
    \item \textbf{Collecte de données :} Conception d'instruments XLSForm avec logiques de saut et contraintes de validation sur \textbf{KoboToolbox}, \textbf{ODK}, CommCare et mWater ; déploiement d'enquêtes multi-sites et géospatiales ; administration des systèmes de collecte.
    \item \textbf{Bases de données et systèmes d'information :} Administration de bases relationnelles (MySQL, PostgreSQL) recevant les données de collecte ; bases propres à chaque projet ; nettoyage, contrôles de cohérence et dédoublonnage des listes de bénéficiaires.
    \item \textbf{Qualité des données :} Protocoles \textbf{DQA} établis et supervisés ; vérification par remontée des chiffres rapportés jusqu'aux registres sources ; suivi des actions correctives jusqu'à clôture.
    \item \textbf{Analyse et restitution :} Analyse statistique et échantillonnage (économiste appliqué de formation) ; \textbf{Excel (avancé)}, \textbf{R}, SQL et Python ; tableaux, graphiques et tableaux de bord \textbf{Power BI}, Looker Studio et Quarto pour le suivi périodique des indicateurs et le rapportage bailleur.
    \item \textbf{Cadres MEAL :} Cadres logiques, chaînes de résultats, plans MEAL et tableaux de suivi des indicateurs, avec valeurs de référence, cibles, désagrégations et sources de vérification.
    \item \textbf{Formation et appui aux équipes :} Formation d'agents de collecte et d'équipes partenaires aux outils numériques et aux protocoles de qualité ; guides techniques qui les rendent autonomes.
\end{itemize}}

\section{Expériences MEAL}

\cventry{Oct 2021 -- Août 2024}{Officier de Suivi \& Évaluation}{Impact Youth Project, Caris Foundation International}{Haïti}{}{
\begin{itemize}
    \item \textbf{Collecte et suivi multi-sites :} Conduite du cycle de suivi d'un projet de santé et de VIH auprès de jeunes, sur des cliniques multi-sites : standardisation des formulaires et des flux de données entre sites, consolidation des remontées dans les délais.
    \item \textbf{Encadrement des agents de collecte :} Direction technique des agents de collecte sur l'ensemble des sites --- répartition du travail, contrôle qualité, formation aux procédures --- et référent technique du personnel de suivi-évaluation sur les contrôles de cohérence.
    \item \textbf{Vérification et assurance qualité (DQA) :} Établissement et supervision des protocoles de Data Quality Assessment sur les données intégrées MySQL, avec remontée des chiffres rapportés jusqu'aux registres sources et suivi des actions correctives jusqu'à clôture.
    \item \textbf{Systèmes d'information :} Mise en place et administration des systèmes de collecte numérique et de la base alimentant le rapportage (CommCare, HIVHaiti, MySQL).
    \item \textbf{Rapportage bailleur :} Production et soumission des indicateurs MER destinés à l'USAID via DATIM, la plateforme du PEPFAR bâtie sur DHIS2.
    \item \textbf{Analyse \& restitution :} Rapports d'activité automatisés en Python et SQL, réduisant le temps de traitement manuel de 40 \% ; tableaux de bord programme sur AWS (Quarto) et \textbf{Power BI} pour le suivi périodique des indicateurs.
\end{itemize}}

\cventry{Oct 2023 -- Jan 2026}{Consultant MEAL \& Analyste de Données}{HANWASH}{Haïti / Télétravail}{}{
\begin{itemize}
    \item Raffinement des plans MEAL et des tableaux de suivi des indicateurs alignés sur les normes JMP (OMS/UNICEF) pour des programmes d'eau, d'hygiène et d'assainissement.
    \item Déploiement d'enquêtes géospatiales complexes sur mWater et administration de la base ; tableaux de bord automatisés (mWater, \textbf{Power BI}) connectés aux API de collecte, guides techniques et formation des équipes à leur usage autonome.
\end{itemize}}

\cventry{Jan 2026 -- Mars 2026}{Analyste de Données \& Consultant MEAL}{Anseye Pou Ayiti}{Haïti}{}{
\begin{itemize}
    \item Définition de la stratégie de collecte et élaboration des indicateurs de performance de trois programmes ; codage des instruments XLSForm pour \textbf{ODK} avec logiques de saut et contraintes de validation, administration du système ODK, et \textbf{formation des agents de collecte} à leur utilisation.
\end{itemize}}

\cventry{Nov 2015 -- Août 2023}{Économiste --- Direction de Planification Économique et Sociale (DPES)}{Ministère de la Planification et de la Coopération Externe (MPCE)}{Port-au-Prince}{}{Suivi de l'exécution du Plan Triennal d'Investissement 2014--2016 : indicateurs de projets publics sur l'ensemble des secteurs et des départements géographiques ; unité statistique et modélisation, avec participation à la construction du << Modèle de Planification du Développement >>.}

\section{Expériences en Ingénierie de Données et Systèmes}

\cventry{Févr 2026 -- Présent}{Ingénieur Logiciel (Fullstack)}{Tekkod LLC}{Télétravail}{}{Modernisation d'infrastructure mobile avec fiabilité hors ligne et sécurité des données ; tableaux de bord Looker Studio.}

\cventry{Mars 2021 -- Mai 2024}{Développeur Backend \& Ingénieur de Données}{Tekkod LLC}{Télétravail}{}{Administration de bases relationnelles en accès concurrent ; pipelines ETL (Apache Beam) migrant les données MongoDB vers BigQuery ; services REST sécurisés (FastAPI, JWT) avec contrôle d'accès par rôle.}

\cventry{Nov 2015 -- Mars 2021}{Consultant Logiciel \& Spécialiste Intégration de Données}{CassionSoft (Projet UNOPS)}{Haïti}{}{Intégration de données sur un projet sous contrat des Nations Unies, avec exposition aux procédures et standards documentaires onusiens.}

\section{Formation}
\cvitem{2010--2014}{\textbf{Études Supérieures en Économie Appliquée} --- C.T.P.E.A (Centre de Techniques de Planification et d'Économie Appliquée), Port-au-Prince : cycle de quatre ans en économie et statistique appliquées, au-delà du niveau Bac+2 requis. \textit{Scolarité complétée ; Diplôme d'Études Supérieures en attente de la soutenance du mémoire de sortie --- attestation officielle jointe au dossier.}}
\cvitem{2009--2010}{Baccalauréat (Mention) --- Institution Saint-Louis de Gonzague, Port-au-Prince}
\cvitem{Autoformation}{FlatIron School (Software Engineering), CodeWithMosh, Udemy}

\section{Compétences Techniques}
\cvitem{Collecte}{\textbf{KoboToolbox}, \textbf{ODK (Open Data Kit)}, \textbf{XLSForm}, CommCare, mWater, HIVHaiti, DHIS2/DATIM \textit{(soumission de données)}, protocoles DQA, ITT}
\cvitem{Bureautique \& viz}{\textbf{MS Excel (avancé)}, Word, PowerPoint, GSuite, \textbf{Power BI (expert)}, Looker Studio, Teams et Zoom}
\cvitem{Analyse statistique}{\textbf{R} (RMarkdown/Quarto), SQL, Python (Pandas), modélisation statistique, échantillonnage, analyse qualitative et quantitative}
\cvitem{Bases de données}{MySQL, PostgreSQL, MongoDB, SQLite, Azure, AWS, Google Cloud (BigQuery), Docker, Git}
\cvitem{Langues}{\textbf{Français (maternel)}, \textbf{Créole haïtien (maternel)}, \textbf{Anglais (courant, rédaction de rapports)}}

\section{Références}
\cvitem{}{Trois références professionnelles, avec leurs coordonnées, figurent sur la page jointe au dossier.}

\end{document}
